% Created 2012-12-06 Thu 23:33
\documentclass[upright, contnum]{umemoria}
\usepackage[latin1]{inputenc}
\inputencoding{latin1}
%\usepackage[T1]{fontenc}
%\usepackage{fontspec}
\usepackage{graphicx}
%\defaultfontfeatures{Mapping=tex-text}
%\setmainfont{Linux Libertine O}
%\setmonofont[Scale=0.8]{DejaVu Sans Mono}

\let\evensidemargin\oddsidemargin
\reversemarginpar

\pagestyle{empty}



\title{informe}
\author{Mauricio}
\date{06 December 2012}

\begin{document}

\maketitle

\setcounter{tocdepth}{3}
\tableofcontents
\vspace*{1cm}

\chapter{Introducción}
\label{sec-1}


  Al igual que en el buffet de un restaurante, por mucho que se quisieran
  comer todos los platos, es imposible comer todo lo que uno
  quisiera por razones obvias. Una posibilidad es probar un poco de cada
  comida, para así saber qué es lo más delicioso y comer hasta
  hartarse.
  
  Pero, ¿qué hacer si hay demasiados platos y no se conocen todos? de
  alguna manera hay que saber cuáles hay que probar, si el objetivo es
  comer lo mejor posible. Un amigo, u otras personas que hayan comido
  en el restaurant pueden recomendar una u otra comida,
  lo cual puede servir para orientarse. Entonces se pueden escoger
  muestras de acuerdo a las recomendaciones.

  
  Pasando a un contexto diferente, supóngase que este gran buffet es la
  Web y los distintos platos corresponden a contenido publicado en
  ella. Por lo tanto, dada la gran cantidad de información disponible,
  se hace necesario poder encontrar lo más atractivo de acuerdo a la
  preferencia de los usuarios. Nótese que se está
  haciendo otra suposición importante con esta analogía, y es que se
  está considerando que la información es íntegramente para ser
  \emph{consumida}, y no, por ejemplo, para generar más contenido,
  conocimiento, o para ser utilizada por máquinas, etc. Dentro de
  este contexto se plantea la pregunta de cómo seleccionar el contenido
  más atractivo dentro de todo lo que hay disponible en un momento dado.

  Para hacerse la idea de cuál es la mejor oferta gastronómica de un
  restaurant, sin tener que probar todos los platos, se puede
  preguntar a la gente que come habitualmente ahí qué platos fueron
  los que más les han gustado, y así definir la calidad de éste y/o
  escoger los mejores platos. O bien, filtrar los documentos que
  hablan de un mismo evento de acuerdo a las preferencias de los
  usuarios, y luego presentarlo en la forma de un \emph{resumen} de cada
  uno.

  Este trabajo consistió en el desarrollo de una metodología
  que permite generar resúmenes automáticos de eventos determinados
  a partir de los documentos Web que hablan de éstos. No se hace
  ninguna suposición sobre el \emph{contenido} de estos documentos (pueden
  ser de texto, imágenes, videos, \emph{multimedios} en general). Los
  documentos del resumen son seleccionados de acuerdo a 
  \emph{indicadores sociales}: los elementos con mayor impacto en las redes
  sociales son considerados más importantes.
  
  La metodología diseñada se implementó sobre dos tipos de eventos:
  noticias y conciertos musicales. Para obtenerlos, se utilizó
  el servicio de Google News\footnote{\href{http://news.google.com/}{http://news.google.com/} } y
  Last.fm\footnote{\href{http://last.fm/}{http://last.fm/} }. Los documentos y sus indicadores
  sociales se obtuvieron de la red social
  Twitter\footnote{\href{http://twitter.com/}{http://twitter.com/} }; y utilizando una representación
  adecuada de los documentos como vectores, se utilizó una estrategia
  de clustering para identificar los subtópicos de cada evento y así
  generar un resumen, basándose en los indicadores obtenidos.

  La estructura de este informe es como sigue: en este capítulo se
  comenta el contexto dentro del cual se desarrolló este sistema, las
  contribuciones realizadas, los objetivos y una descripción general
  de la solución; en el siguiente capítulo se discute el estado del
  arte y el marco teórico del cual se desprende este trabajo; el
  capítulo 3 describe más en detalle el problema a resolver, su
  relevancia y sus dificultades, para luego, en el
  capítulo 4, describir la solución implementada, más un par de casos
  de estudio sobre los resultados obtenidos, para finalizar con las
  conclusiones de este trabajo en el capítulo 5.

\section{Contexto y Motivación}
\label{sec-1.1}

   
   La tasa de crecimiento de la cantidad de datos en la Web, y en
   particular, en las \emph{redes sociales online} (OSN, \emph{Online Social Networks}),
   es de tal magnitud que se vuelve necesario encontrar formas de
   filtrar y buscar sólo la información relevante dentro de todas las
   fuentes que hablan del mismo evento. 

   Por otra parte, para poder entender lo que está pasando en el
   mundo a partir de los puntos de vista de toda las personas que
   manifiestan su opinión o los hechos recabados en las redes
   sociales, se hace necesario poder procesar toda esta
   información. Nuevamente, dado el gran volumen de datos, es
   necesario poder resumir esta información para ser rápidamente
   consumida. 

   En el contexto de las redes sociales online, cada día los usuarios
   publican  millones de mensajes cortos con respecto a distintos
   tópicos, ya sean conversacionales, sobre eventos noticiosos, 
   etc.\footnote{Pear Analytics. Twitter Study \href{http://es.scribd.com/doc/18548460/Pear-Analytics-Twitter-Study-August-2009}{http://es.scribd.com/doc/18548460/Pear-Analytics-Twitter-Study-August-2009} }
   Además, el auge de los teléfonos inteligentes o \emph{smartphones} con mayor
   capacidad de procesamiento e integrados con todo tipo de sensores
   (cámaras fotográficas, de vídeo, acelerómetro, osciloscopio, etc.),
   hace posible el generar aun más información y
   e incluso en tiempo real sobre lo que acontece en el mundo, en
   Internet, o bien sobre el estado particular de cada usuario.

   Este aumento y evolución de la generación de datos no sólo influye en la
   riqueza de éstos, sino también en el comportamiento de los usuarios
   a lo largo del tiempo. Actualmente,  una gran parte de éstos valora
   más el contenido de tipo multimedia (imágenes y videos)
   en las redes sociales online\footnote{The Rise of Visual Social Media \href{http://www.fastcompany.com/3000794/rise-visual-social-media}{http://www.fastcompany.com/3000794/rise-visual-social-media}. En el artículo se menciona un estudio sobre comportamiendo y preferencias de los usuarios en las redes sociales llevado a cabo por ROI Research: \href{http://www.slideshare.net/performics_us/performics-life-on-demand-2012-summary-deck}{http://www.slideshare.net/performics\_us/performics-life-on-demand-2012-summary-deck} }. 
   Este contenido no sólo corresponde a documentos Web en el sentido
   tradicional, sino a objetos de más alto nivel como \emph{tweets}
   (mensajes cortos de la red social Twitter), imágenes (de servicios
   como Instagram\footnote{\href{http://instagr.am/}{http://instagr.am/} },
   Tumblr\footnote{\href{http://tumblr.com/}{http://tumblr.com/} }, etc.), vídeos
   (Youtube\footnote{\href{http://youtube.com/}{http://youtube.com/} }, Vimeo\footnote{\href{http://vimeo.com/}{http://vimeo.com/} }) e
   incluso sonidos (Soundcloud\footnote{\href{http://soundcloud.com/}{http://soundcloud.com/} }).
   Se hace entonces necesario encontrar formas para satisfacer estas
   necesidades de los usuarios, las cuales ya han sido
   abordadas en parte, como la generación de
   resúmenes automáticos orientado a motores de búsqueda, o la
   determinación de la relevancia tanto de documentos en la Web como de
   mensajes en las redes sociales.

   Surge como motivación el poder identificar y extraer contenido
   relevante de lo que está pasando en las redes sociales,
   clasificados dentro del concepto de eventos, y además avanzar un
   paso más arriba en el nivel de abstracción: considerar los
   documentos no por su contenido textual, lo que permite abarcar
   imágenes, vídeos, sonidos y multimedios en general. Algunas
   aplicaciones directas de esto son, entre otras:

\begin{itemize}
\item Ayudar al trabajo periodístico mediante una colección de
     contenido multimedia relacionado a un evento noticioso. Por
     ejemplo, la versión online de Radio
     Biobío\footnote{\href{http://www.biobiochile.cl/}{http://www.biobiochile.cl/} } frecuentemente publica
     breves artículos sobre sucesos que tienen impacto en las redes
     sociales, mostrando un pequeño conjunto de mensajes con
     comentarios de la gente\footnote{Como muestra: \href{http://www.biobiochile.cl/2012/12/01/aporte-de-lustrabotas-de-santiago-a-la-teleton-provoca-admiracion-en-redes-sociales.shtml}{http://www.biobiochile.cl/2012/12/01/aporte-de-lustrabotas-de-santiago-a-la-teleton-provoca-admiracion-en-redes-sociales.shtml}, y \href{http://www.biobiochile.cl/2012/12/01/rechazo-provocan-condicionamientos-de-compra-de-ripley-y-unimarc-para-donar-a-la-teleton.shtml}{http://www.biobiochile.cl/2012/12/01/rechazo-provocan-condicionamientos-de-compra-de-ripley-y-unimarc-para-donar-a-la-teleton.shtml} }. 
     Una aplicación directa involucraría
     considerar además contenido multimedia, y organizar este
     contenido de acuerdo a la relevancia que tiene dentro de las
     redes.
\item Enriquecer la búsqueda en la Web a través de contenido
     multimedia. Una persona buscando información sobre un concierto
     podría obtener imágenes y vídeos de éste fácilmente una vez
     identificado el concierto.
\item Siguiendo lo anterior, un grupo musical podría obtener toda la
     información multimedia asociada a su concierto, tanto para sus
     fans como para ellos mismos, potenciando su popularidad.
\item Poder distinguir entre eventos similares rápidamente. Por
     ejemplo, un usuario que desee obtener información sobre ``Gaza'',
     puede referirse tanto a la banda de música como al conflicto en
     Israel. El poder distinguir rápidamente mediante una imagen o un
     vídeo acelera mucho el proceso. \emph{Una imagen vale más que mil palabras}.
\end{itemize}
   El sistema implementado es una primera aproximación que puede
   satisfacer los ejemplos mencionados.

\section{Objetivos}
\label{sec-1.2}

\subsection{Objetivo general}
\label{sec-1.2.1}


    El objetivo principal de este trabajo fue el siguiente:

    Diseñar e implementar un sistema que permita generar resúmenes
    automáticos de \emph{eventos}: información temporal publicada en redes
    sociales online sobre sucesos en particular. Esta información se
    basa en contenido textual y multimedial generado por los usuarios
    de estas redes, cuya relevancia se basa en el impacto generado en
    éstas.


\subsection{Objetivos específicos}
\label{sec-1.2.2}


\begin{enumerate}
\item Extraer datos relacionados a eventos en la Web, principalmente
       aquellos generados en redes sociales online. Estos datos pueden
       componerse tanto de información textual como multimedial.
\item Agrupar la información extraída de un evento en subtópicos.
\item Seleccionar los elementos más relevantes de cada subtópico para
       producir un resumen del evento.
\item Analizar la efectividad de la metodología propuesta sobre un
       conjunto de eventos noticiosos y conciertos.
\end{enumerate}
\section{Descripción general de la solución}
\label{sec-1.3}

   
   Para enfrentar el problema se diseñó una solución que consiste en
   tres componentes para resolverlo, resultando cada una en diferente
   grado de complejidad, siendo posible además ser mejoradas en el
   futuro. Para probar la viabilidad de la metodología, ésta fue
   aplicada sobre un conjunto de casos de prueba.

   En particular:

\begin{itemize}
\item Se llevó a cabo una metodología para la obtención de documentos y
     enriquecerlos con datos obtenidos de fuentes sociales;
\item Se diseñó un procedimiento que separar estos documentos en
     \emph{clusters}, \emph{sin considerar su contenido}. Sólo se utilizó la
     información social asociada; y
\item Se diseñó además un procedimiento para \emph{rankear} u ordenar los
     resultados de acuerdo a \emph{relevancia}, siendo ésta medida de
     acuerdo a la información social asociada a los documentos
     generados.
\end{itemize}
   Las componentes diseñadas fueron las siguientes:

\begin{enumerate}
\item La que obtiene descripciones de eventos a partir de fuentes de
      éstos en la Web, enriqueciéndolos con información social;
\item Otra componente que procesa y separa los documentos a partir de
      la información social; genera \emph{objetos Web} y los separa en
      subtópicos de cada evento, respectivamente; y
\item La componente que entrega los $k$ documentos más relevantes por
      cada evento obtenido, basándose en los subtópicos identificados.
\end{enumerate}
   Se utilizaron las API de Google News como de Last.fm para la
   obtención de eventos: noticias y conciertos, respectivamente. Para
   el enriquecimiento de los eventos se utilizó la información social
   que provee Twitter y su API de búsqueda de \emph{tweets}. De la misma
   forma, se consideraron los metadatos de los mismos mensajes para medir
   la relevancia de los documentos generados. 

   Un documento es identificado por la URL que lo ubica en la Web. El
   contenido no es más que la concatenación de los tweets que
   mencionan al documento. Se realizó una limpieza y preprocesamiento
   de los datos, quitando las \emph{stopwords} y realizando \emph{stemming}
   sobre el contenido en texto. Luego, se aplicó \emph{tf-idf} sobre los
   documentos, representándolos como vectores en el \emph{space vector    model}. Para identificar los subtópicos de un evento se utilizó el 
   algoritmo de clustering $k$-means sobre los vectores. 

   Para el ranking de los documentos se decidió usar una ponderación
   simple sobre una serie de indicadores que dependen de los tweets y
   de las URLs de cada evento.

   Entre las herramientas utilizadas, se usó lenguaje de
   programación Python, varias librerías para el manejo de datos
   (tales como \texttt{nltk}, \texttt{scipy}, \texttt{scikit-learn}, por nombrar las más
   importantes), el sistema de almacenamiento Redis, entre otras
   herramientas que son mencionadas en la descripción detallada de la
   solución.

\chapter{Antecedentes}
\label{sec-2}

  
  Para poder describir correctamente tanto el problema como la
  solución implementada, es necesario dar los punteros y conceptos
  básicos que los involucran. En este capítulo se discutirán los
  siguientes tópicos:

\begin{itemize}
\item La red social Twitter, la cual es utilizada como fuente de datos y
    documentos para este trabajo.
\item Clustering de documentos, y en general, estrategias de clustering,
    las cuales tienen muchas aplicaciones prácticas. En este trabajo
    fue utilizada una de estas estrategias para poder determinar los
    subtópicos de cada evento.
\item La identificación automática de eventos, la cual, si bien se
    utilizó un enfoque más simple para este trabajo, sirve para
    indicar en qué aspectos es posible extender este trabajo en el
    futuro.
\item Resúmenes automáticos: una sucinta definición, y algunos enfoques
    que han existido en el tiempo para este procedimiento.
\item Ranking de documentos, o cómo generar órdenes de acuerdo a
    relevancia.
\end{itemize}
  Casi todos estos tópicos, a excepción del primero, involucran
  técnicas de Minería de Datos, Recuperación de la Información y
  Aprendizaje de Máquinas, entre otras áreas.

\section{Twitter}
\label{sec-2.1}

   Twitter es una red social online que permite conectar a
   personas mediante la comunicación de mensajes cortos, rápidos y   frecuentes\footnote{\href{https://support.twitter.com/groups/31-twitter-basics/topics/104-welcome-to-twitter-support/articles/13920-get-to-know-twitter-new-user-faq}{https://support.twitter.com/groups/31-twitter-basics/topics/104-welcome-to-twitter-support/articles/13920-get-to-know-twitter-new-user-faq} }. Estos
   mensajes son publicados en el perfil del usuario que los emite,
   pueden ser vistos directamente por los seguidores de este usuario o
   ser vistos directamente en el perfil o buscándolos mediante una
   funcionalidad que provee el servicio. Además, un usuario puede
   \emph{seguir} a otros para poder ver en su \emph{timeline} o perfil privado 
   los mensajes de todos a quienes sigue.
   
\begin{figure}[h!b]
  \centering
  \includegraphics[width=12cm]{./img/twitter.png}
  \caption[Timeline de Twitter.]
   {Timeline de Twitter. En éste se ve una lista en orden cronológico
  de los tweets generados por los usuarios que sigue el usuario
  actual. Además, el sitio incluye tweets promocionados por cuentas
  que pagan por dicho servicio, como es el caso del primer tweet en la
  lista.}
\end{figure}
   
   Estos mensajes, o \emph{tweets}, sólo son cadenas de caracteres con
   metadatos que el mismo servicio asigna una vez enviado a la red
   social. Desde sus inicios (año 2007) se han añadido algunas capacidades
   adicionales a estos mensajes, como la de poner URLs, imágenes,
   vídeos, etc. Además, existen varias convenciones que han surgido a
   lo largo del tiempo. A continuación se describe una lista de tipos
   de mensajes que existen en Twitter, originados por estas convenciones:

\begin{enumerate}
\item Respuestas o \emph{replies}: son mensajes del tipo \texttt{@usuario [texto]},
      que ocurren usualmente en una conversación entre dos usuarios.
\item Menciones o \emph{mentions}: un poco más general a una respuesta, el
      nombre del usuario mencionado puede estar en cualquier parte del
      mensaje. La diferencia semántica es que no se le habla
      ``directamente'' al usuario mencionado, como en una respuesta, sino
      que sólo es mencionado por si el mensaje es de su interés o no.
\item \emph{Retweets}: son mensajes del tipo \texttt{RT @usuario: [texto]}. Ocurren
      cuando se quiere compartir el mensaje de otro usuario, o citarlo
      para mencionarlo en el mismo mensaje.
\item \emph{Hashtags}: son palabras precedidas por el caracter \#, que indican
      un identificador a cierto evento o suceso dentro o fuera de la
      red. Suelen usarse para categorizar de cierta forma un tópico, pero
      son libres de usarse como los usuarios quieran.
\item Mensaje simple: un mensaje sin menciones ni hashtags.
\end{enumerate}
  Ejemplos:

\begin{itemize}
\item Mensaje simple: \texttt{Jason Funk disipa patitos};
\item Respuesta: \texttt{@jason estoy de acuerdo con lo que dices};
\item Mención: \texttt{creo que @jason es una cumbre de sabiduría};
\item Retweet: \texttt{RT @jason: Jason Funk disipa patitos}; y
\item Hashtag: \texttt{Estoy escribiendo mi memoria \#dcc \#summarization}
\end{itemize}
  Estos mensajes están limitados a 140 caracteres de extensión. Sumando
  esto a la integración de la red con otros servicios y dispositivos, y
  a la cantidad de mensajes publicados cada minuto, permite utilizar
  esta red como una gran fuente de datos.

  Twitter además provee varios servicios adicionales, como por ejemplo,
  un servicio de acortamiento de URLs, para permitir incluir una URL
  larga sin perjudicar la cantidad de caracteres restantes para el
  mensaje; un servicio de alojamiento de fotos y vídeos, para hacer más
  sencilla la publicación de mensajes multimedia desde dispositivos
  móviles; un servicio de búsqueda que permite buscar una cantidad
  determinada de tweets sobre un término de búsqueda o un hashtag,
  entre otros servicios.

   

   
   
   

\section{Clustering de documentos}
\label{sec-2.2}


   El análisis de clusters o clustering es el proceso de encontrar
   grupos de objetos, tal que los objetos en un grupo sean similares
   entre sí (o relacionados) y que sean diferentes (o no relacionados)
   a los objetos de otros grupos. Algunas aplicaciones del análisis de
   clusters son, entre otras:
\begin{itemize}
\item Encontrar clusters naturales y describir sus propiedades (\emph{data understanding});
\item Encontrar agrupamientos útiles (\emph{data class identification});
\item Encontrar representantes de grupos homogéneos (\emph{data reduction});
\item Encontrar perturbaciones aleatorias de los datos (\emph{noise detection});
\item Encontrar objetos inusuales (\emph{outliers detection});
\item etc.
\end{itemize}
   Se denomina cluster a un grupo de objetos, mientras que
   clustering puede referirse al conjunto de clusters o al proceso de
   encontrarlos. Existen diversos tipos de procesos de clustering, una
   de las distinciones más importantes es entre los clusters
   jerárquicos y los particionales:
\begin{itemize}
\item Un clustering jerárquico es un conjunto de clusters anidados,
     organizados más bien como un árbol. Cortando el árbol en
     cualquier nivel da como resultado un clustering potencialmente
     distinto.
\item El clustering particional es un conjunto de clusters de forma de
     partición del conjunto total, es decir, cada objeto está
     contenido en un sólo subconjunto o cluster.
\end{itemize}
   Para describir el proceso aplicado a documentos, primero se
   describirán los modelos de representación más importantes para,
   de forma de definir la noción de documento y luego los
   algoritmos de clustering aplicados a éstos.

\subsection{Modelos de representación de documentos}
\label{sec-2.2.1}


    \subsubsection{Standard Boolean Model}

    El modelo booleano es un modelo de representación de
    documentos. En él, los documentos son vectores de \emph{términos}:

    $$d = (w_1,w_2,\ldots,w_m)$$

    Donde un término es un $n$-grama del texto del documento.

    \begin{defn} Un $n$-grama es una secuencia contigua de $n$ ítems a
    partir de un texto. \end{defn}

    La definición de ítem dependerá de la aplicación: en lenguaje
    natural el texto a su vez dependerá del idioma, por ejemplo, si el
    texto está en inglés o en japonés, la distinción entre palabras
    es distinta para cada uno. 

    No existe una medida de ``similitud'' como tal en este modelo, sino
    que se considera el calce exacto entre los términos de una query
    $q$ y un documento $d$. La query puede ser una consulta hecha por
    un usuario al conjunto de documentos, o bien un documento del
    mismo conjunto.

    Una consulta es una fórmula de lógica proposicional que pide los
    documentos que contengan o no ciertos términos.

    \subsubsection{Bag of words Model}

    En el modelo Bag of Words un documento $d$ es representado como un
    conjunto de pares $(w_i, f_i)$, $i\in[1..m']$, donde $w_i$ es un
    término del documento, $f_i$ es la frecuencia de $w_i$ en el
    mismo, y $m'$ es la cantidad de términos distintos en el
    documento. 

    La ventaja principal por sobre el modelo anterior es que permite
    hacer calces parciales entre consultas y documentos. Este modelo
    es comúnmente utilizado para hacer clasificación de documentos,
    por ejemplo, para determinar si un correo electrónico es o no
    spam.

    \subsubsection{Vector Space Model}

    El \emph{Vector Space Model} es un modelo un poco más general que el
    anterior. Un documento $d$ es representado como un vector de pesos
    asociados a los términos:

    $$d = (f(w_1), f(w_2), \ldots, f(w_m))$$

    Cada dimensión de este vector corresponde al peso asociado a un
    término del documento. 

    El peso puede ser directamente la frecuencia del término dentro
    del documento:

    $$\freq(w,d) = |\{w : w \in d\}|$$

    O bien, normalizar esta frecuencia para evitar que documentos más
    largos sean más relevantes sólo por su extensión:

    $$\tf_0(w,d) = \left\{ 
    \begin{array}{l l}
    1 & \quad \textrm{si $w \in d$}\\
    0 & \quad \textrm{si no}\\
    \end{array} \right.$$

    $\tf_0$ o \emph{Term Frequency} es una primera aproximación a medir la
    frecuencia de un término en un documento. Sin embargo, esta nueva
    aproximación sufre de la desventaja de que ahora un documento con
    una ocurrencia del término será igual de relevante que algún
    documento que mencione varias veces el término (por ejemplo, un
    diccionario que tiene el término una vez contra un artículo sobre
    el tema). Otra alternativa, considera no castigar demasiado a los
    documentos con pocas ocurrencias, pero tampoco beneficiar mucho a
    los que tengan muchas:

    $$\tf_1(w,d) = 1 + \log(\freq(w,d))$$

    La solución más utilizada considera la proporción con respecto al
    término con más ocurrencias, para esto, se normaliza por el tamaño
    del documento:

    $$\tf(w,d) = \frac{\freq(w,d)}{max\{\freq(t,d) : t \in d\}}$$

    Otro problema que tiene utilizar esta medida como los pesos de los
    términos, es que un término muy repetido entre todos los
    documentos que hablan de un mismo tema puede significar que no es
    muy relevante (por ejemplo, las \emph{stopwords} o palabras vacías, son
    por lo general las preposiciones, artículos, pronombres,
    etc.). Para esto, se considera además ponderar por el inverso de
    la frecuencia entre los documentos; es decir, un término frecuente
    entre todos los documentos ve su peso castigado a diferencia de un
    término que sólo es mencionado una vez en un documento. Esta
    medida es llamada \emph{Inverse Document Frequency} o $\idf$:

    $$\idf(t, D) = \log \frac{ |D| } {1 + |\{d \in D : t \in d\}| }$$

    Finalmente, el peso de un término es la ponderación de su
    frecuencia dentro del documento con el inverso de la frecuencia
    entre los documentos, o $\tfidf$:

    $$\tfidf(t,d,D) = \tf(t,d) \times \idf(t,D)$$

    Para el resto de este trabajo se considerará la representación de
    documentos en este modelo.

\subsection{Medidas de similitud}
\label{sec-2.2.2}

      
    A lo largo de los años, dos formas han sido las usuales para
    determinar similitud entre
    documentos\cite{Zhao02criterionfunctions}: calculando el coseno de
    los dos vectores, y la distancia euclidiana.

    El coseno de dos documentos $d_i$ y $d_j$ se define como

    $$\cos(d_i, d_j) = \frac{d_i^td_j}{\|d_i\|\|d_j\|}$$

    Si ambos $d_i$ y $d_j$ son iguales, entonces la fórmula anterior
    se evalúa a 1, y 0 si no tienen nada en común, es decir, ambos
    vectores son ortogonales.

    La distancia euclidiana entre $d_i$ y $d_j$ se define como

    $$d(d_i,d_j) = \sqrt{(d_i-d_j)^t(d_i-d_j)} = \|d_i-d_j\|$$

    En este caso, si la distancia es 0, entonces ambos documentos son
    idénticos. Y si ambos son ortogonales, la distancia será
    $\sqrt{2}$. 

    Si ambos documentos están normalizados (su norma es 1), entonces
    ambas medidas serán muy parecidas, aun así siendo una de similitud
    y la otra de distancia.

\subsection{Algoritmos de clustering}
\label{sec-2.2.3}


    Se considerará la siguiente definición para el problema de
    clustering:
    
    Dado un entero $k$ y un conjunto de $n$ puntos en $\mathbb{R}^v$,
    el objetivo es determinar $k$ puntos tal que se minimice $\phi$,
    la suma de las distancias al cuadrado de cada punto y su centro
    más cercano.

    Este problema es NP-hard, sin embargo, existe un algoritmo que
    permite realizar una búsqueda local y determinar $k$ centros 
    en un tiempo razonable\footnote{Lloyd, S.; , ``Least squares quantization in PCM,'' Information Theory, IEEE Transactions on , vol.28, no.2, pp. 129- 137, Mar 1982 }.
    
    K-means es un algoritmo de clustering particional, en el cual cada
    cluster tiene un \emph{centroide} asociado, esto es, típicamente, un punto
    el cual es el promedio de todos los puntos del cluster, y no
    necesariamente corresponde a un dato real.
    
    El algoritmo genera $k$ clusters, donde $k$ es un parámetro del
    algoritmo, tal que cada punto pertence al cluster cuyo centroide es el
    más cercano a éste.
    
\begin{algorithm}[H]
 \KwData{$k>0$, conjunto $D$ de puntos}
 \KwResult{Asignación de cada punto a un cluster $i \in [1..k]$}
 Seleccionar $k$ puntos como centroides iniciales\;
 \While{Los centroides cambien a cada iteración}{
  Formar $k$ clusters asignando todos los puntos al centroide más cercano\;
  Recalcular los centroides de cada cluster\;
 }
 \caption{K-means}
\end{algorithm}

    Los clusters generados dependerán de la elección inicial de los
    centroides, y usualmente basta con pocas iteraciones para la
    convergencia. La complejidad de este algoritmo es $O(nkIv)$, donde $n$
    es el número de puntos, $k$ el parámetro de la cantidad de clusters,
    $I$ es la cantidad de iteraciones que hará el algoritmo y $v$ es la
    cantidad de dimensiones de los vectores.

    K-means también puede verse desde el enfoque de optimizar una
    función criterio. La medida más común es la 
    \emph{suma del error cuadrado} o Sum of Squared Error (SSE). Para cada
    punto, su \emph{error} es la distancia al cluster más cercano:
    $$SSE = \sum_{i=1}^k\sum_{x\in C_i} \dist^2(m_i,x)$$
    Donde $x$ es un punto en el cluster $C_i$, y $m_i$ es su
    centroide. Dados dos clusters, en cada paso se escoge el que tenga
    menor error.
    Un problema importante de K-means es escoger los centroides
    iniciales. Dado que la probabilidad de acertar a los centroides
    ``reales'' (precisión) es muy baja, es posible que un clustering
    tenga un mal error cuadrático, de hecho, el radio competitivo con
    la solución óptima es no acotado incluso para $k$ y $n$ fijos. 
    Existen varias formas de evitar
    este problema, por ejemplo, ejecutando varias veces K-means, 
    usar clustering jerárquico para determinar los centroides
    iniciales, determinar más de $k$ centroides y luego elegir los $k$
    más separados, etc.

    Un enfoque utilizado es el de usar el algoritmo
    K-means++\cite{Arthur:2007:KAC:1283383.1283494}, el cual tiene
    como objetivo determinar los centroides iniciales para
    K-Means. Este algoritmo garantiza un radio competitivo de 
    $O(\log k)$ con respecto al error cuadrático esperado del óptimo.
    El algoritmo K-means++ escoge los centroides con cierta
    probabilidad que depende de la distancia al centroide más cercano,
    lo cual garantiza una buena elección inicial, y luego continúa con
    el algoritmo usual. Sea $D(x)$ la distancia de $x$ al centroide
    más cercano. El algoritmo es como sigue:

\begin{algorithm}[H]
 \KwData{$k>0$, conjunto $D$ de puntos}
 \KwResult{Asignación de cada punto a un cluster $i \in [1..k]$}
Escoger un centroide $c_1$ al azar uniformemente de $D$\;
Escoger un nuevo centroide $c_i$, escogiendo $x \in D$ con probabilidad $\frac{D(x)^2}{\sum_{x' \in D}D(x')^2}$\;
Repetir paso 1 hasta haber escogido $k$ centros\;
Proceder con K-means estándar\;
 \caption{K-means++}
\end{algorithm}
    

\subsection{Evaluación de clusters}
\label{sec-2.2.4}


    Existen dos enfoques para evaluar clusterings: validación interna
    y externa\footnote{\href{http://en.wikipedia.org/wiki/Cluster\_analysis\#Evaluation\_of\_clustering\_results}{http://en.wikipedia.org/wiki/Cluster\_analysis\#Evaluation\_of\_clustering\_results} } La validación interna considera evaluar el clustering
    con respecto a los datos que han sido agrupados (es decir, sin
    información adicional), mientras que la evaluación externa
    considera datos adicionales, como datos pre-clasificados (conjunto
    de prueba, de entrenamiento, etc), o los clusters ``reales''.

    Por otra parte, dos de las métricas más usadas para evaluar
    clusterings, segun \cite{Zhao02criterionfunctions} son:

\begin{itemize}
\item \textbf{Entropía}: mide cómo son distribuidas las clases de documentos
      entre cada cluster. Se define formalmente como
      
      $$E(C_r) = -\frac{1}{\log q}\sum_{i=1}^q\frac{n^i_r}{n_r}\log\frac{n^i_r}{n_r}$$

      donde $q$ es el número de clases en el dataset, y $n^i_r$ es el
      número de documentos de la $i$-ésima clase que fue asignado al
      $r$-ésimo cluster $C_r$. La entropía del clustering se define
      como la suma ponderada de la entropía de cada cluster:

      $$E(\mathfrak{C}) = \sum_{r=1}^k \frac{n_r}{n} E(C_r)$$

      Un clustering perfecto tendrá clusters tal que cada cluster
      contenga documentos de una sola clase, en ese caso la entropía
      será 0. En general, conviene tener bajos valores de entropía.
\item \textbf{Pureza}: mide de qué forma cada cluster contiene documentos de
      una clase. La pureza de un cluster $C_r$ se define como

      $$P(C_r) = \frac{1}{n_r} max_i\{n^i_r\}$$

      la cual es la fracción del tamaño del cluster\ldots{}.\ldots{}.\ldots{}.\ldots{}.\ldots{}.\ldots{}.\ldots{}.\ldots{}.
\end{itemize}
\section{Identificación automática de eventos}
\label{sec-2.3}


   La identificación automática de eventos consiste en, dado un
   conjunto de documentos, donde cada documento está asociado a un
   evento (desconocido), poder particionar el conjunto de
   documentos en clusters, de forma que cada cluster corresponda a
   todos los documentos asociados a un evento.

   Se considerará la definición de ``evento'' dada por \cite{Yang:1999:LAD:630307.630471}. 

   \begin{defn} Un \emph{evento} es un suceso que ocurre en un período de tiempo
   determinado y en un lugar específico. \end{defn}

   El poder identificar eventos a partir de documentos publicados en
   los medios sociales permite mejorar la navegación de estos eventos,
   al mejorar la búsqueda tanto local como de motores de búsqueda.

   

\section{Resúmenes automáticos}
\label{sec-2.4}

\subsection{Evaluación de resúmenes}
\label{sec-2.4.1}

\section{Ranking de documentos}
\label{sec-2.5}

   
\chapter{Especificación del Problema}
\label{sec-3}


  El problema a resolver consiste en poder \emph{resumir} eventos en base
  al contenido textual y multimedial de la información publicada en
  medios sociales, tales como las redes sociales online.

  Se considerará un evento como una ocurrencia en el mundo real
  con un período de tiempo asociado y un conjunto de mensajes, de
  volumen considerable, que discutan la ocurrencia y además publicados
  dentro del período de tiempo. 

  El resumir un evento consiste en entregar un subconjunto de
  mensajes o de documentos mencionados en ellos que tengan
  relación directa con el evento en cuestión. Por ejemplo, el evento
  \texttt{Anef anuncia movilización nacional} contiene muchos mensajes, tales
  como:
\begin{itemize}
\item \texttt{RT @econtingencia: Reajuste: Trabajadores del sector público}
    \texttt{convocan a movilizaciones: La negociación entra a su recta final:}
    \texttt{El l... h ...}
\item \texttt{CUT de La Araucanía se adherirá a marcha nacional en rechazo a}
    \texttt{reajuste salarial: La Central Unitaria de}
    \texttt{Trabaja... http://t.co/IcgohlTR}
\item \texttt{Reajuste: Trabajadores del sector público convocan a}
    \texttt{movilizaciones http://t.co/pWC7q6NI}
\item \texttt{RT @biobio: CUT de La Araucanía se adherirá a marcha nacional en}
    \texttt{rechazo a reajuste salarial ofrecido por el Gobierno}
    \texttt{http://t.co/IOieDido}
\end{itemize}
  El resumen consistirá en una colección de objetos, los cuales pueden
  ser o bien mensajes o bien algunos de los documentos mencionados en
  ellos (en este caso representados por las URLs contenidas en los
  mensajes), de forma que esta colección sea de tamaño mucho menor que
  el total de mensajes, y además que los elementos de esta colección
  representen de alguna forma todos los puntos de vista o aspectos
  importantes del evento.

  Se deben tomar en consideración las características de estos
  mensajes: muchos de ellos tendrán errores de ortografía, redacción,
  gramática, o semántica (que hablen de otro tema). Esto se traduce
  como el ruido de los datos. Una solución a este problema debe tomar
  en consideración el ruido intrínseco de los datos. Además, estos
  mensajes son breves y no expresan en profunidad un punto de vista
  del evento, sino que usualmente pueden tener una o dos
  características cruciales que permitan distinguirlos de entre todos
  los que existen, por ejemplo, una mención a una imagen o a un
  vídeo que explique una parte del evento.


\chapter{Descripción de la Solución}
\label{sec-4}


  La solución propuesta consistió en generar un resumen utilizando
  documentos representados apropiadamente con la información de medios
  sociales que los mencionan. La relevancia de estos documentos
  también se obtuvo con la información social más otros indicadores de
  la solución parcial. 

  Para esto, se asume que existe una fuente de eventos, de esta forma,
  el enfoque de la solución radica principalmente en la generación de
  los resúmenes y no en la identificación de eventos en medios
  sociales. Sin embargo, para la implementación sí fue necesario
  enfrentar el problema de identificar eventos, usando un enfoque
  práctico para obtenerlos.

  La figura FIGURA muestra de manera general el modelo propuesto.

  FIGURA MODELO

  DESCRIPCION DEL MODELO

  En la Sección \ref{descdet} se describe el modelo utilizado para la
  implementación, con una especificación detallada de la solución. En la
  sección \ref{impl} se describe la metodología de desarrollo y la
  implementación práctica realizada para representar el modelo
  formal, la cual considera la obtención de datos dentro del
  proceso. Luego se comentan los problemas técnicos que fueron
  enfrentados para terminar en la sección \ref{casosest}, donde se discuten
  un par de casos de estudio sobre los resultados obtenidos.

\section{Descripción detallada}
\label{sec-4.1}

\label{descdet}


   Para enfrentar el problema descrito se decidió utilizar una
   representación apropiada de los documentos que permita abstraerse
   de su contenido, utilizando la información social asociada a
   éstos. Se decidió que el problema de identificación de eventos está
   fuera del alcance de este trabajo, por lo cual, se asume que los
   eventos son dados como input a la solución diseñada. A partir de un
   evento determinado, se identifican los  subtópicos del evento
   utilizando los documentos, y luego, para cada 
   subtópico se determinan los documentos más relevantes utilizando
   esta información social.

   Para esto, fue necesario contar con dos \emph{fuentes de datos}: una
   fuente de eventos y otra de \emph{contenido social}, en la forma de
   mensajes y actualizaciones de estado.  

   Se asumió que estas fuentes satisfacen los siguientes
   requerimientos:

\begin{itemize}
\item La fuente de eventos debe entregar una lista de eventos
     rápidamente, la cual debe contener los siguientes datos para
     cada entrada:

\begin{itemize}
\item Un título del evento y \emph{términos asociados}. Los términos
       asociados son breves frases o palabras que describan al
       evento, como por ejemplo, tags o etiquetas.
\item Como datos opcionales: breve descripción del evento, fecha de
       inicio y término, ubicación y direcciones Web.
\end{itemize}

\item La o las fuentes de contenido social deben entregar una lista de
     mensajes, con algunos metadatos tales como la fecha de creación,
     si el mensaje fue compartido, etc. Además, algunos datos sobre el
     autor del mensaje, como la cantidad de conexiones en la red, y
     en general, datos que permitan comparar dos autores.
\end{itemize}
   Utilizando estas dos fuentes, el siguiente paso luego de obtener
   una lista de eventos fue enriquecerlos utilizando las fuentes
   sociales, generando documentos del tipo 
   $d = (s_1, s_2, \ldots, s_m)$, donde  $s_i$, $i \in [1..m]$ 
   es un mensaje de alguna fuente social, con los
   metadatos asociados. El documento es identificado por la URI de
   algún documento en la Web, de forma que todos los mensajes que
   contengan una URI en particular, corresponderán al mismo documento
   $d$.

   A continuación, utilizando alguna representación adecuada 
   (\emph{vector space model}, \emph{bag of words}, etc.), se generaron clusters
   de documentos de un mismo evento, identificando los subtópicos. Con
   ellos fue posible generar un resumen que abarcara todos los
   aspectos del evento, en contraste con seleccionar directamente los
   documentos más relevantes del evento en su conjunto, lo cual puede
   dejar puntos de vista sin ser considerados por su extensión. 

   Finalmente se seleccionaron los $k>0$ documentos más representativos de
   cada cluster, utilizando como criterio los metadatos de los
   mensajes de la fuente social. De esta forma, se ordenan los
   documentos dejando como más ``relevantes'' los que más interés atrae
   de los usuarios.

\section{Metodología de desarrollo e implementación}
\label{sec-4.2}

\label{impl}

   La implementación consistió en las siguientes etapas:

\begin{itemize}
\item Obtención del dataset de eventos y documentos.

     Para esto, se utilizaron dos fuentes de eventos: Google News y
     Last.fm para recolectar noticias y conciertos musicales
     (incluidos festivales), respectivamente. Como fuente de datos
     sociales se utilizó la red social Twitter, que dispone de una API
     para realizar búsquedas por \emph{keywords}. Esta etapa comprendió
     la recolección de eventos y de documentos con información social
     asociada a éstos.
\item Generación de clusters para la identificación de subtópicos para
     cada evento.

     Una vez identificados los eventos y los documentos asociados, se
     generaron clusters usando el algoritmo K-means con K-means++ para
     la inicialización. Se impuso un valor de $k=5$ clusters por
     evento. En la sección de casos de estudio se discuten las
     alternativas y la evaluación de algunos clusters del dataset.
\item Extracción de documentos relevantes para cada evento.

     Una vez identificados los subtópicos de cada cluster, se
     extrajeron de éstos los documentos más relevantes, utilizando la
     información social de cada uno de ellos, en conjunto con otros
     indicadores globales del clustering (como que incluyan una URL
     dentro de las más mencionadas dentro del cluster, entre
     otras). Estos documentos corresponden al output o salida del
     sistema.
\end{itemize}
\subsection{Obtención de datos}
\label{sec-4.2.1}

    Se describe a continuación el proceso diseñado para la obtención de
    datos, tanto de eventos como de documentos con sus indicadores
    sociales respectivos.

    Las etapas de generación del dataset son las siguientes:

\begin{itemize}
\item Recolección de eventos (noticias y conciertos);
\item Enriquecimiento de los eventos existentes mediante tweets; e
\item Identificación de documentos a partir de los tweets por cada evento.

      Se recolectaron datos (eventos y tweets) desde el 19 de noviembre de
      2012 hasta el 30 de noviembre del mismo año, todos los días
      desde la medianoche hasta que el proceso termina exitosamente.
\end{itemize}
\subsubsection{Recolección de eventos}

Se consideraron dos tipos de eventos para el sistema: noticias y
conciertos musicales. Los conciertos incluyen festivales de varios
artistas.

\begin{itemize}
\item Noticias
  Para obtener las noticias, se utilizó el servicio de Google
  News\footnote{\href{http://news.google.com}{http://news.google.com} }. Existe una API (en proceso de
  obsolescencia\footnote{\href{http://googlecode.blogspot.com/2011/05/spring-cleaning-for-some-of-our-apis.html}{http://googlecode.blogspot.com/2011/05/spring-cleaning-for-some-of-our-apis.html} }, 
  pero funcional a la fecha de este trabajo) que permite
  obtener no sólo los titulares y breve descripción de cada noticia,
  sino también un conjunto de entre 4-10 noticias relacionadas de otras
  fuentes. Esto sirvió para alimentar los términos de búsqueda para la
  etapa siguiente. Se guardaron los siguientes datos de una noticia:

\begin{itemize}
\item Título,
\item Descripción,
\item URL de la fuente, y
\item Titulares de las noticias relacionadas.
\end{itemize}

\item Conciertos
  Utilizando el servicio de Last.fm para obtener los conciertos y
  festivales de una ubicación en
  particular\footnote{\href{http://www.lastfm.es/api/show/geo.getEvents}{http://www.lastfm.es/api/show/geo.getEvents} }, se
  obtuvieron los conciertos y festivales de las siguientes
  ubicaciones:

\begin{itemize}
\item Santiago, Chile;
\item Londres, Inglaterra;
\item Glastonbury, Inglaterra;
\item Las Vegas, Nevada, EE.UU.; y
\item Estocolmo, Suecia.
\end{itemize}

\item Título del evento (concierto o festival);
\item Artistas que participan; y
\item Fechas de inicio y término (esta última no siempre está como
    dato).

  Además de otros datos descriptivos, como la ubicación, descripción
  breve, sitio web de la banda o festival, etc.
\end{itemize}
Cada vez que se obtienen los eventos se vuelven a obtener los
conciertos, pero sólo agregando los nuevos. Las noticias siempre son
nuevas, aun así por implementación no se consideraron los repetidos.
  
\subsubsection{Enriquecimiento de eventos}

Se obtuvieron tweets utilizando el servicio de búsqueda que provee
Twitter en su
API\footnote{\href{https://dev.twitter.com/docs/api/1.1/get/search/tweets}{https://dev.twitter.com/docs/api/1.1/get/search/tweets} }. El
objetivo es enriquecer los eventos con la información social que hay
en la Web sobre éstos. 

Para cada uno de los eventos obtenidos en la fase anterior, se
utilizaron los términos de búsqueda asociados a ellos: los titulares
de las noticias relacionadas y los nombres de los artistas para los
eventos noticiosos y musicales, respectivamente.

\begin{itemize}
\item Para las noticias, se hace una búsqueda en Twitter de los titulares
  al mismo tiempo en que se obtienen de Google News, y nuevamente al
  día siguiente, es decir, 2 búsquedas por cada titular de un evento.
  Se quitan las tildes y caracteres no \texttt{ASCII} y las stopwords, para
  evitar problemas con la implementación y no hacer calce de stopwords
  en la búsqueda de Twitter.
\item Para los conciertos y festivales, se utilizaron los nombres de los
  artistas y del evento como términos de búsqueda. De acuerdo a la
  información asociada al evento, se busca por una mayor cantidad de
  días:

\begin{itemize}
\item Se busca desde un día antes de inicio del evento;
\item Si está presente la fecha de término del evento, se busca cada día
    dentro del intervalo ``fecha de inicio'' a ``fecha de término'' hasta
    tres días terminado el evento.
\item Si no está presente la fecha de término (por ejemplo, un concierto
    o un festival de un día), se busca hasta tres días pasada la fecha
    de inicio.
\end{itemize}

\end{itemize}
\subsubsection{Identificación de documentos a partir de tweets}

    Luego de obtener los tweets asociados a cada evento, el siguiente
    paso fue generar los documentos que fueron usados para la
    generación de los resúmenes. Nuevamente, el modelo consistió en que cada
    documento se modeló como un vector de palabras, donde el
    identificador del documento es una URL, y sus componentes
    corresponden al contenido de los tweets que tienen esa URL en el
    texto del mensaje.

    El caso en el que un tweet no tenía ninguna URL en su contenido
    fue abordado de la siguiente forma: la URL asociada es una tal que
    representa al mismo tweet (utilizando el servicio de Twitter), y
    el contenido de ese documento es el mismo tweet, de forma de no
    dejar el tweet sin ser representado.

    Este proceso fue abordado recorriendo todos los eventos del
    dataset, observando todos los tweets asociados a cada evento,
    extrayendo la URL si es que hay alguna y guardando el documento
    con el nuevo tweet. Se marcan los tweets observados para no tener
    que repetir el proceso, ya que es intensivo en conexión a la red.

    Dada la condición breve de los mensajes publicados en la red
    social, muchos de los usuarios y/o servicios que publican mensajes
    con una URL en su interior suelen utilizar \emph{acortadores} 
    (o \emph{url shorteners}) para los enlaces, y así no utilizar mucho
    espacio dentro de un
    mensaje. Otra ventaja que ofrecen es que algunos servicios como
    Bit.ly\footnote{\href{http://bit.ly}{http://bit.ly} } dan estadísticas sobre los visitantes a
    estos enlaces (y así saber quiénes vienen de cierta red social u
    otra, por ejemplo). Twitter, a su vez, actualmente también ofrece
    acortamiento de URLs por defecto. Esto suele producir que un enlace
    acortado se resuelva a otro enlace también acortado, por lo que es
    necesario resolver la URL completa para evitar duplicados o
    \emph{pseudo-duplicados} (en el caso en que dos URLs sintácticamente
    distintas apunten al mismo recurso). EN LA FIGURA\ldots{}...

    FIGURA DE LINKS CORTOS

    Por lo anterior, una vez identificada la URL del texto de un
    tweet, se resuelve su URL completa (que puede ya serlo de
    antemano), lo que consume recursos de ancho de banda y
    tiempo. 

\subsection{Identificación de subtópicos}
\label{sec-4.2.2}


    Una vez recolectados tanto los eventos como generados los
    documentos asociados, se procedió a identificar los subtópicos de
    cada evento. Para esto, se utilizó el algoritmo K-means para
    construir clusters a partir de todos los documentos de un solo
    evento.

    Como se mencionó anteriormente, los documentos consisten en
    vectores del tipo $d_{\textrm{URL}}=(t_1,t_2,\ldots,t_m)$, donde
    $t_i$ es el tweet $i$-ésimo que contiene a URL dentro del texto
    del mensaje. Para poder aplicar un algoritmo de clustering, fue
    necesario procesar nuevamente estos documentos para representarlos
    como vectores usando el vector space model. El procedimiento
    consta de dos partes, ``normalizar'' los documentos, limpiando los
    términos que puedan afectar al clustering, y luego aplicar
    $\tfidf$ sobre los documentos normalizados:

\begin{algorithm}[H]
 \KwData{Conjunto de documentos $D_e$}
 \KwResult{Conjunto de strings $D'_e$, }
 $D'_e \leftarrow \emptyset$\; 
 \For{documento $d \in D_e$}{
   $d' \leftarrow \varepsilon$\;
   \For{tweet $t \in d$}{
   $t' \leftarrow \texttt{clean}(t)$\;
   $d' \leftarrow \texttt{concat}(d',t')$\;
   }
   $D'_e \leftarrow D'_e \cup \{d'\}$\;
 }
 \caption{Preprocesamiento de documentos}
\end{algorithm}
    
    Donde $\varepsilon$ es el string vacío, \texttt{concat}$(a,b)$ retorna
    la concatenación de $a$ y $b$, y \texttt{clean} realiza las siguientes
    operaciones sobre el tweet:

\begin{itemize}
\item Remueve las URLs que contenga el texto;
\item Remueve todas las menciones;
\item Quita los caracteres ``\#'', dejando los \emph{hashtags} intactos;
\item Quita las tildes, acentos, stopwords; y
\item Realiza stemming en español o inglés dependiendo del idioma del
      evento dado por los metadatos de éste. Se utilizó el Snowball
      Stemmer\footnote{\href{http://snowball.tartarus.org/}{http://snowball.tartarus.org/} } para esto.
\end{itemize}
    Una vez normalizados los documentos, se convierten a la
    representación como vectores con pesos por cada término:

\begin{algorithm}[H]
\KwData{Conjunto de strings $D'_e$, vocabulario $V$ de palabras de
$D'_e$}
\KwResult{Conjunto de vectores $D''_e$, representados en vector space
model}
$D''_e \leftarrow \emptyset$ \;
\For{documento $d \in D'_e$} {
   $d' \leftarrow \texttt{map}(\tfidf(\cdotp, d, D'_e), \texttt{words}(d))$\;
}
$D''_e \leftarrow D''_e \cup \{d'\}$\;
\caption{Transformación de documentos a vector space model}
\end{algorithm}

    Donde \texttt{map}$(f,l)$ mapea la
    función o procedimiento $f$ a cada elemento de la lista $l$,
    retornando una nueva lista $l'$ en la cual a cada elemento se le
    aplicó $f$ y \texttt{words} retorna una lista con las palabras de
    $d'$. Este procedimiento retorna un conjunto de vectores en
    $\tfidf$, lo que permite entonces aplicar un algoritmo de
    clustering para identificar subtópicos.

    Se utilizó el algoritmo Mini Batch K-means, el CUAL\ldots{} COMPLETAR
    ACA Y MODIFICAR ANTES DONDE DIGA KMEANS!!!!!. Al aplicarlo sobre
    el conjunto de vectores, éste retorna una \emph{lista de etiquetas}
    $L$, de largo $n$, la cantidad de documentos utilizados. Cada
    $L_i \in [0..k-1], i \in [1..n]$ indica a cuál cluster corresponde
    el $i$-ésimo documento, lo que permite fácilmente filtrar y
    recuperar los clusters por separado, para pasar a la siguiente
    etapa.

\subsection{Ranking de documentos}
\label{sec-4.2.3}

    


\section{Desafíos técnicos}
\label{sec-4.3}

\subsection{Restricciones de la API de Twitter}
\label{sec-4.3.1}


   La API de búsqueda de Twitter permite obtener tweets de acuerdo a un
   término de búsqueda. Se utilizó este servicio para enriquecer los
   eventos con información social utilizando como términos de búsqueda
   tanto los títulos de las noticias como los nombres de los artistas
   para las noticias y los conciertos, respectivamente. 
   
   Funciona de la siguiente forma: cada vez que se hace un request a la
   URL dada por el servicio, éste retorna a lo más 100 tweets por página, con un
   máximo de 15 páginas (indicando en el request qué página queremos
   consultar), dando como total hasta 1500 tweets por búsqueda. Existirán
   términos de búsqueda que no presenten ningún resultado  (ya sea por
   estar mal escritos o simplemente que no sean un tópico de discusión), o por
   el contrario, que se generen más tweets que los retornados por la
   búsqueda por cada ventana de tiempo que demore ésta (por ejemplo, un
   \emph{trending topic} o tópico que sea muy mencionado en la red social).
   
   Existe una limitación de uso de este servicio: sólo es posible hacer
   hasta 180 requests por cada 15 minutos, o 1 request cada 5
   segundos. Además, sólo retorna tweets de hasta 7 días de antigüedad, y
   sus resultados no son necesariamente en tiempo real y su estabilidad
   varía de acuerdo a factores externos.
   
   Los tweets retornados vienen en formato \texttt{JSON} (\emph{Javascript Simple Object Notation}),
   e incluyen varios metadatos sobre el tweet aparte de los principales,
   como autor, fecha, contenido. Algunos de estos metadatos son:
   
\begin{itemize}
\item Cantidad de \emph{retweets} hechos hasta la fecha;
\item Si posee alguna URL o \emph{hashtag} en el texto;
\item Si es una \emph{mención} a otro usuario;
\item La ubicación de donde se envió el tweet;
\item etc.
\end{itemize}
  Además, incluye datos sobre el autor, como por ejemplo:

\begin{itemize}
\item Si la cuenta está \emph{verificada};
\item La cantidad de seguidores del usuario;
\item Cantidad de amigos (seguidores que también lo siguen);
\item Cantidad de tweets;
\item Su descripción, y si incluye alguna URL, etc;
\item Ubicación (dada por el mismo usuario);
\item Fecha de creación de la cuenta;
\item etc.
\end{itemize}
\section{Casos de estudio}
\label{sec-4.4}

\label{casosest}

\chapter{Conclusiones}
\label{sec-5}

\section{Resumen del trabajo realizado}
\label{sec-5.1}

\section{Objetivos alcanzados}
\label{sec-5.2}

\section{Relevancia del trabajo realizado}
\label{sec-5.3}

\section{Trabajo futuro}
\label{sec-5.4}


   


























\end{document}
